% ==================== 第三章：问题三（价格反馈与动态滚动优化） ====================
\section{价格反馈与动态滚动优化策略}
\label{sec:model3}

本章聚焦在价格—供给存在相互作用、信息逐年滚动更新的背景下，如何以“年”为阶段实施动态规划：每一年度先利用上一年度的真实产量、售价与销量反馈来更新参数分布，再在面向未来若干情景的基础上进行含下行风险控制的再优化，从而输出下一年度的地块—季节—作物配置。与前两章相比，这里更强调价格弹性与销量上限的双重约束对收益分布尾部的影响，并通过引入 CVaR（条件在险价值）在目标函数中显式惩罚“坏情景”的收益，进而提升保底表现。

为便于理解，我们先以既有结果中的基线收入曲线回顾价格—供给反馈的形状与拐点特征，进而给出动态滚动的建模框架与算法流程；随后展示一组代表性数值实验与敏感性分析，以说明滚动优化在高波动市场中的稳定性优势。


\begin{figure}[H]
  \centering
  \includegraphics[width=0.82\linewidth]{model3_revenue_curve.pdf}
  \caption{基线收入—供给曲线示意（用于刻画价格弹性与容量上限对总收入的影响）}
  \label{fig:m3-revenue-curve}
\end{figure}

\subsection{模型动机与设定}
现实市场中，作物的年度均价并非外生常数，供给增加通常压低价格，同时还受到渠道容量与销量上限的约束。基于上一年度的“计划—实际—成交”三元组，我们将价格与销量分别建模为关于“有效供给”的分段线性或光滑可微（如对数线性）的函数族，再辅以不确定项以反映气象、品质与市场情绪带来的扰动。这样得到的价格—供给反馈函数在小供给区近似线性增长，而在接近容量上限时边际收益显著递减，对应到收益曲线便呈现明显的凹性与拐点（见图~\ref{fig:m3-revenue-curve}）。

为便于后续求解与分析，我们给出两个常用的价格—供给函数原型：
\begin{align}
  &\text{线性价格：}\qquad p(q) = a - bq,\quad a>0,\,b>0,\\
  &\text{对数线性：}\qquad \ln p(q) = \ln p_0 - \eta q\;\Rightarrow\; p(q)=p_0 e^{-\eta q},\quad \eta>0.
\end{align}
考虑销量上限（或渠道容量）$C$ 与有效销量 $s(q)=\min\{q,\,C\}$，则年度总收入可表述为
\begin{equation}
  R(q) = p(q)\,s(q) = p(q)\,\min\{q,\,C\},
\end{equation}
其一阶导数在 $q=C$ 处出现折点：在 $q<C$ 区间，$\frac{\mathrm{d}R}{\mathrm{d}q}=p'(q)q + p(q)$；在 $q>C$ 区间，$\frac{\mathrm{d}R}{\mathrm{d}q}=p'(q)C$。这正对应图~\ref{fig:m3-revenue-curve} 的拐点与“容量饱和”现象。若进一步计入单位变动成本 $c$ 与固定成本 $F$，则利润为 $\Pi(q)=R(q)-cq-F$，在滚动优化中我们以 $\Pi$ 的分布为核心关切对象。

年度决策向量为“年—季—地块—作物”的面积矩阵 $x_{t,\,s,\,i,\,k}$，并配合二进制变量 $z_{t,\,s,\,i,\,k}$ 刻画互斥与重茬约束。标注$y_{t,\,s,\,i,\,k}$ 为产量，$u_{t,\,k}$ 为当年售出量，则有效供给 $q_t=\sum_{s,i,k}y_{t,\,s,\,i,\,k}$，销量受 $u_t\le C_t$ 与 $u_t\le q_t$ 共同约束，价格 $p_t$ 由 $q_t$ 反馈决定。

\subsubsection{变量与约束总览}
核心变量包括：$x$（面积，连续）、$z$（作物-地块-季节是否种植，二进制）、$q$（总供给，连续）、$u$（销量，连续）、$p$（价格，连续）、$r$（收入，连续）。关键约束示意：
\begin{align}
  &\text{地块面积：}& \sum_k x_{t,s,i,k} &\le A_{s,i}, & x_{t,s,i,k} &\ge 0, \\
  &\text{互斥/重茬：}& \sum_k z_{t,s,i,k} &\le 1, & x_{t,s,i,k} &\le M\,z_{t,s,i,k}, \\
  &\text{产量定义：}& y_{t,s,i,k} &= \theta_{t,s,i,k}\,x_{t,s,i,k}, \\
  &\text{总供给：}& q_t &= \sum_{s,i,k} y_{t,s,i,k}, \\
  &\text{销量上限：}& 0\le u_t &\le \min\{q_t,\,C_t\}, \\
  &\text{价格反馈：}& p_t &= f(q_t;\,\Theta_t), \\
  &\text{收入与利润：}& r_t &= p_t\,u_t, & \Pi_t &= r_t - c^\top x_t - F_t.
\end{align}
其中 $\theta_{t,s,i,k}$ 可视为不确定单产参数；$f(\cdot)$ 为价格—供给反馈函数，后续采用分段线性近似以进入 MIP。

与设定配套的年度决策向量仍是“年—季—地块—作物”的种植面积变量，二进制变量用于刻画重茬限制与地块—季节的互斥关系。重要差异在于目标函数的随机性：年度收益取决于当年实现的单产、价格与销量，其联合分布在每年开始时据历史数据和外部信息进行更新，并据此生成未来多情景以供优化。

\begin{figure}[H]
  \centering
  \includegraphics[width=0.78\linewidth]{model3_sales_vs_capacity.pdf}
  \caption{销量—产能的配比关系示意（容量约束导致的瓶颈区间）}
  \label{fig:m3-sales-capacity}
\end{figure}

\subsection{价格—供给—销量反馈与识别}
价格函数的识别遵循“解释性—稳定性—可解性”三项权衡：第一，优先采用分段线性或对数线性规格，使其与线性/混合整数约束保持良好耦合；第二，通过稳健回归或分位回归降低异常点影响；第三，将销量上限以硬约束嵌入，避免非现实的高销量—高价格组合。识别完成后，我们以置信区间或盒式不确定集表征参数区间，供下一步的情景生成使用。作为直观参考，图~\ref{fig:m3-sales-capacity} 显示了销量—产能在瓶颈区间的典型“饱和”现象。

在对数线性规格下，可采用带先验收缩的稳健回归：
\begin{equation}
  \ln p_t = \beta_0 + \beta_1 q_t + \beta_2 \mathbb{1}\{q_t\!>\!C_t\} + \varepsilon_t,\qquad \varepsilon_t\sim\text{Laplace}(0,\,\sigma),
\end{equation}
或使用分位回归估计 $\beta(\tau)$ 来捕捉尾部敏感性。容量 $C_t$ 可由“渠道小时数×处理能力”或“合同规模”给出，其年度更新可用随机游走或回归于设施投资、合作方数量等解释变量。

\subsubsection{数据清洗与稳健性}
对历史 $\{(q_t, p_t, u_t, C_t)\}$ 需执行：去极值/截尾、节令与质量分组校正、渠道波动哑变量控制、分箱后再回归以降低非线性影响。推荐使用 Huber 或 Tukey 损失；对 $\ln p$ 回归的残差以 Ljung–Box 检验序列相关，并用 Newey–West 协方差修正置信区间。

\begin{figure}[H]
  \centering
  \includegraphics[width=0.8\linewidth]{model3_profit_distribution_baseline.pdf}
  \caption{确定性方案在历史波动下的收益分布（作为滚动优化的对照基线）}
  \label{fig:m3-profit-dist-baseline}
\end{figure}

\subsection{情景生成与参数更新}
每年初，我们根据过去一年的实际产量、成交价与销量，对单产、价格弹性、容量上限与成本等参数的先验分布进行贝叶斯式更新或滑动窗口估计，随后在更新后的分布上进行蒙特卡洛采样或拉丁超立方抽样，形成代表性情景集。为控制样本效率，可按“价格-销量-单产”的相关结构进行协方差一致化与主成分降维。为保证下行覆盖，额外在尾部分位（如 5\%）采样加密，便于 CVaR 估计稳定。

在正态—正态共轭下，设 $\theta$ 为某线性价格模型参数向量，先验 $\theta\sim\mathcal N(\mu_0,\,\Sigma_0)$，新样本 $\{(q_t,\,p_t)\}$ 满足 $p_t=X_t\theta+\varepsilon_t,\ \varepsilon_t\sim\mathcal N(0,\,\sigma^2 I)$，则后验为
\begin{align}
  \Sigma_t^{-1}&=\Sigma_{t-1}^{-1}+\tfrac{1}{\sigma^2}X_t^\top X_t, &
  \mu_t&=\Sigma_t\Big(\Sigma_{t-1}^{-1}\mu_{t-1}+\tfrac{1}{\sigma^2}X_t^\top p_t\Big).
\end{align}
若使用稳健误差（如 Laplace），可通过 EM 或变分法得到近似后验。情景生成时，我们先按相关结构对 $(\theta,\,C,\,Y)$ 协同采样，再通过 $p(q)$ 与 $s(q)$ 传播至收益与利润的分布。

\subsubsection{相关结构与采样}
设联合随机向量 $\xi=(\theta, C, Y, \epsilon_p)$ 具有经验协方差 $\hat\Sigma$，先做主成分分解 $\hat\Sigma=V\Lambda V^\top$ 并保留 $m$ 个主成分，随后在 $\mathcal N(0, I_m)$ 上做拉丁超立方抽样，最后回到原空间：$\xi = \hat\mu + V_m\Lambda_m^{1/2} z$。尾部加密通过在目标方向（如最小利润方向）上加权采样实现。

图~\ref{fig:m3-profit-dist-baseline} 给出了前两章基线方案在历史波动下的收益分布，作为后续滚动 + CVaR 策略的对照。图~\ref{fig:m3-spider-base} 与图~\ref{fig:m3-heatmap-switch} 则分别刻画了基线方案对关键参数的敏感度和地块—作物切换的空间分布，为设计滚动启发式提供直观依据。

% （合并至后文并排子图）原“敏感性（基线）”单图已合并到并排图中，标签保留在子图中

\begin{figure}[H]
  \centering
  \includegraphics[width=0.78\linewidth]{model3_switch_heatmap.pdf}
  \caption{地块—作物切换热图（利于识别滚动优化中需重点稳定的子区）}
  \label{fig:m3-heatmap-switch}
\end{figure}

\subsection{含 CVaR 的两阶段滚动优化模型}
为在年度滚动框架下控制收益尾部风险，我们引入 CVaR 于目标函数。具体地，年度目标可表述为“最大化期望收益减去 \(\lambda\) 倍 CVaR@\(\alpha\)”，其中 \(\lambda\) 表示风险厌恶强度，\(\alpha\) 为置信水平（如 95\%）。在样本近似下，CVaR 可通过引入辅助变量与线性松弛约束得到线性/混合整数可解的等价形式；当与价格—供给分段线性反馈函数合并时，整体仍保持 MIP 结构，与前章求解器无缝衔接。

采用 Rockafellar–Uryasev 公式，设情景 $s=1,\dots,S$ 的年度损失为 $\ell_s=-\Pi_s$，则
\begin{align}
  \mathrm{CVaR}_\alpha(\ell) &= \min_{\zeta}\Big\{\zeta + \frac{1}{(1-\alpha)S}\sum_{s=1}^S [\ell_s-\zeta]_+\Big\}\\
  &\equiv \min_{\zeta, u_s\ge 0}\Big\{\zeta + \frac{1}{(1-\alpha)S}\sum_{s=1}^S u_s:\ u_s\ge \ell_s-\zeta\Big\}.
\end{align}
在 MIP 中引入 $\zeta,\,u_s$ 及样本权重，即可线性化地刻画 CVaR 项。年度目标为
\begin{equation}
  \max_{x, z, \zeta, u}\ \mathbb{E}[\Pi_s]\; -\; \lambda\,\mathrm{CVaR}_\alpha(\ell)
\end{equation}
约束：作物—地块—季节、容量与价格—供给反馈等。
切换惩罚可通过相邻年度解的差分范数惩罚来实现，例如 $\kappa\,\|x^{(t)}-x^{(t-1)}\|_1$，或在变量层面以额外二进制“切换”指示器保障。

\subsubsection{价格分段线性化与 MIP}
对 $p(q)$ 采用断点集 $\{b_j\}_{j=0}^J$ 与段斜率 $\{\phi_j\}$，以 SOS2 或凸组合变量 $\lambda_j$ 实现：$q=\sum_j b_j\lambda_j,\ p=\sum_j p(b_j)\lambda_j,\ \sum_j\lambda_j=1,\ \lambda_j\ge 0$。在容量邻域加密断点以提升拟合质量。若需对 $r=p\,u$ 的双线性进行线性化，可用 McCormick 包络在 $[\underline p,\overline p]\times[0,C]$ 上近似。

为展示与鲁棒预算 \(\Gamma\) 的关系，我们在图~\ref{fig:m3-gamma-frontier} 与图~\ref{fig:m3-regret-curve} 中并列给出“鲁棒预算—收益前沿”与“后悔值—风险尾部”的对照：CVaR 以概率加权的方式约束尾部，而鲁棒优化在不确定集内采取最坏情景的确定性保证。二者在稳健性目标上一致，但在保守度与可解释性上各有优势。

\begin{figure}[H]
  \centering
  \begin{subfigure}[t]{0.48\linewidth}
    \centering
    \includegraphics[width=\linewidth]{model3_gamma_frontier.pdf}
    \caption{鲁棒预算 \(\Gamma\)—收益前沿（与 CVaR 前沿的概念性对照）}
    \label{fig:m3-gamma-frontier}
  \end{subfigure}\hfill
  \begin{subfigure}[t]{0.48\linewidth}
    \centering
    \includegraphics[width=\linewidth]{model3_regret_curve.pdf}
    \caption{后悔值曲线（展示不同策略在最不利情景下的收益让渡）}
    \label{fig:m3-regret-curve}
  \end{subfigure}
  \caption{鲁棒预算与后悔值的并排对比}
\end{figure}

\subsection{求解策略与算法流程}
滚动优化的关键在于“更新—生成—优化—执行—评估”的闭环。为降低年度求解时间与防止过度反复跳变，我们引入地块切换惩罚与分解启发式，并对价格函数在拐点附近做自适应分段加密。整体流程如算法~\ref{alg:rolling-cvar} 所示。

具体而言：
\begin{itemize}
  \item \textbf{分段自适应：}对 $p(q)$ 在 $q\approx C$ 邻域加密断点，以较少段数逼近非线性拐点；其余区间采用较粗分段以节省变量与约束。
  \item \textbf{分解启发式：}先对年度总供给 $q$、销量 $u$ 做连续松弛求近似，再在固定 $(q,u)$ 的子问题中细化到地块—季节层面的离散分配；交替迭代数轮直至收敛或达到时间上限。
  \item \textbf{热启动与惩罚续接：}用上年解 $x^{(t-1)}$ 热启动，并逐步提高切换惩罚 $\kappa$，避免一开始过度保守导致的停滞。
  \item \textbf{并行样本评估：}在候选解之间以并行评估（外层蒙特卡洛）筛选，保留“高期望—低尾部”的非支配解，再做精调。
\end{itemize}

\subsubsection{复杂度与收敛性讨论}
分段点数 $J$ 与情景数 $S$ 共同决定模型规模，MIP 变量量级约为 $\mathcal O(|x|+|z|+J+S)$。在滚动框架下，每年只求一次主问题，计算量与单次全期优化相当或更低。对于分解启发式，可证明在固定切换惩罚与断点集下的交替最优化收敛到坐标极值。

\begin{algorithm}[htbp]
  \caption{滚动优化 + CVaR 决策流程（以年为阶段）}
  \label{alg:rolling-cvar}
  \begin{algorithmic}[1]
    \State 输入：基期方案 $x^{(0)}$、历史数据 $\mathcal{D}_{0}$、风险参数 $\alpha,\,\lambda$
    \For{年度 $t=1,\dots,7$}
      \State 基于 $\mathcal{D}_{t-1}$ 更新参数先验，得到 $\Theta_t$
      \State 在 $\Theta_t$ 上生成情景集 $\{\xi_s\}_{s=1}^S$（含尾部分位加密）
      \State 构建含 CVaR 的两阶段 MIP 并加入切换惩罚与分段线性价格函数
      \State 求解得到当年方案 $x^{(t)}$，执行并观测实际结果 $y^{(t)}$
      \State 更新数据集 $\mathcal{D}_{t} \leftarrow \mathcal{D}_{t-1} \cup \{x^{(t)}, y^{(t)}\}$
    \EndFor
    \State 输出：全期滚动方案与评估指标（期望/分位收益、CVaR、切换频率等）
  \end{algorithmic}
\end{algorithm}

\subsection{数值实验与对比分析}
我们以前两章的确定性与鲁棒解作为对照基线，对滚动 + CVaR 策略进行外层 1000 次蒙特卡洛评估。结果显示：当 $\lambda$ 适中且 $\alpha$ 取 95\% 时，期望收益仅小幅让渡，但最坏 5\% 分位与 CVaR 指标显著改善。图~\ref{fig:m3-struct-share} 呈现了在滚动策略下的作物结构演变；图~\ref{fig:m3-util-robust} 和图~\ref{fig:m3-revenue-capped} 显示了资源利用与价格—供给拐点对整体收益弹性的抑制作用。

\subsubsection{实验设置}
作物集 $\mathcal K$ 含 6 种作物，地块 $\mathcal I$ 含 24 块，季节 $\mathcal S$ 含 3 季；年周期 $T=7$。单产以对数正态波动，价格以对数线性反馈，容量按随机游走更新。风险参数 $(\alpha,\lambda)=(0.95,0.6)$ 为主设定，$\kappa$ 取分段递增序列以实现平滑切换。

\subsubsection{结果解读}
滚动 + CVaR 在“尾部—期望”权衡上优于单纯鲁棒与确定性方案，结构上倾向于在容量瓶颈季节分散高弹性作物，减少对价格敏感的集中投放；资源利用曲线显示其在关键资源上的峰值更平缓，表明对极端情景的更好适配能力。

\begin{figure}[H]
  \centering
  \caption{滚动策略相关图的并排整合展示}
  \begin{subfigure}[t]{0.48\linewidth}
    \centering
    \includegraphics[width=\linewidth]{model3_stack_area.pdf}
    \caption{滚动策略下的作物结构演变（以基线结构图为参照）}
    \label{fig:m3-struct-share}
  \end{subfigure}\hfill
  \begin{subfigure}[t]{0.48\linewidth}
    \centering
    \includegraphics[width=\linewidth]{model3_resource_util_line_robust.pdf}
    \caption{资源利用率在鲁棒/滚动策略下的对比（线图）}
    \label{fig:m3-util-robust}
  \end{subfigure}
  \\[0.6em]
  \begin{subfigure}[t]{0.7\linewidth}
    \centering
    \includegraphics[width=\linewidth]{model3_revenue_curve_capped.pdf}
    \caption{价格—供给约束下的收入曲线与拐点（滚动策略的收益弹性参考）}
    \label{fig:m3-revenue-capped}
  \end{subfigure}
\end{figure}

\subsection{敏感性与稳健性讨论}
为分离“参数不确定性”和“决策切换成本”两类来源的影响，我们分别在价格弹性、销量上限、单产方差与切换惩罚系数上做一维与交互敏感性实验。图~\ref{fig:m3-spider-robust} 汇总了鲁棒与滚动策略下的关键敏感度排序，图~\ref{fig:m3-spider-compare} 则直观对比了“基线 vs 滚动 + CVaR”的敏感性差异；图~\ref{fig:m3-violate} 展示了在不同保守度与风险参数下的约束违约风险对比，验证滚动策略在高波动区间的约束可行性提升。

在实验设定中，取价格弹性 $\eta\in[0.02,0.15]$、容量 $C\in[0.6\bar{q},1.2\bar{q}]$、单产方差 $\sigma_Y^2$ 与切换惩罚 $\kappa\in[0,\kappa_{\max}]$。我们记录对期望利润、5\%分位、$\mathrm{CVaR}_{95\%}$ 的敏感度：
\begin{align}
  S_{\text{exp}}(\eta) &= \frac{\partial \mathbb{E}[\Pi]}{\partial \eta}, \\
  S_{\text{cvar}}(\kappa) &= \frac{\partial \,\mathrm{CVaR}_{0.95}(-\Pi)}{\partial \kappa}, \\
  S_{p5}(C) &= \frac{\partial \,\text{Quantile}_{0.05}(\Pi)}{\partial C}.
\end{align}
结果显示，滚动 + CVaR 相较确定性/鲁棒单策，显著降低了 $\eta$ 与 $C$ 在尾部指标上的敏感度，对切换惩罚呈现“先快后缓”的边际改善。

% 合并：基线 与 滚动 + CVaR 并排子图，扩展加入“对比”与“违约风险”子图
\begin{figure}[H]
  \centering
  \begin{subfigure}[t]{0.48\linewidth}
    \centering
    \includegraphics[width=\linewidth]{model3_sensitivity_spider_baseline.pdf}
    \caption{敏感性（基线）}
    \label{fig:m3-spider-base}
  \end{subfigure}\hfill
  \begin{subfigure}[t]{0.48\linewidth}
    \centering
    \includegraphics[width=\linewidth]{model3_sensitivity_spider_robust.pdf}
    \caption{敏感性（滚动 + CVaR）}
    \label{fig:m3-spider-robust}
  \end{subfigure}
  \\[0.6em]
  \begin{subfigure}[t]{0.48\linewidth}
    \centering
    \includegraphics[width=\linewidth]{model3_sensitivity_spider_compare_side_by_side.pdf}
    \caption{敏感性对比（基线 vs 滚动 + CVaR）}
    \label{fig:m3-spider-compare}
  \end{subfigure}\hfill
  \begin{subfigure}[t]{0.48\linewidth}
    \centering
    \includegraphics[width=\linewidth]{model3_violation_risk_bar.pdf}
    \caption{约束违约风险柱状图（不同保守度与风险参数的对比）}
    \label{fig:m3-violate}
  \end{subfigure}
  \caption{敏感性与风险的并排对比（基线 vs 滚动 + CVaR）}
\end{figure}

\FloatBarrier

\subsubsection{管理启示与实施建议}
滚动优化将“学习—决策—评估”闭环化，使价格—供给反馈被逐年消化并反映到可执行的地块配置中；CVaR 的引入显式压制了收益分布尾部，从而在销量受限与价格大幅波动时提升保底收益。与单次全局确定性规划相比，滚动策略的优势在于可持续纠偏与算力友好；与鲁棒优化相比，滚动+CVaR 的尾部控制更贴近概率意义下的风险管理实践。管理上，建议在年度计划编制中同步上报“尾部指标”（如 CVaR@95\%、最坏分位），并对地块切换设定“红线频率”，以兼顾收益与实施稳定性。

% 额外参考图：作物单产与地块适宜性的年度剖面（与滚动更新相关）与资源热图并排
\begin{figure}[H]
  \centering
  \begin{subfigure}[t]{0.48\linewidth}
    \centering
    \includegraphics[width=\linewidth]{model3_yield_bar.pdf}
    \caption{单产剖面图（为滚动更新中的产量分布提供直观基线）}
    \label{fig:m3-yield-profile}
  \end{subfigure}\hfill
  \begin{subfigure}[t]{0.48\linewidth}
    \centering
    \includegraphics[width=\linewidth]{model3_resource_heatmap.pdf}
    \caption{资源使用热图（识别滚动优化中需优先协调的瓶颈环节）}
    \label{fig:m3-resource-heat}
  \end{subfigure}
  \caption{年度单产与资源使用的并排对比}
\end{figure}

\subsection{本章小结}
本章在价格—供给存在反馈、销量受容量上限约束且信息逐年更新的背景下，提出了以“年”为阶段的滚动优化框架，并在年度目标中引入 CVaR 对尾部风险进行显式惩罚。核心建模要素包括：价格—供给—销量的分段线性化反馈、基于贝叶斯更新与相关一致化的情景生成、Rockafellar–Uryasev\cite{arletti2024directional} 形式的 CVaR 线性化以及地块切换惩罚与启发式分解，整体保持与 MIP 求解器的良好耦合。

主要结论与管理启示如下：
\begin{itemize}
  \item \textbf{可解性与稳定性：}通过在容量拐点附近加密分段、对双线性项采用 McCormick 包络，以及在年度间加入切换惩罚，兼顾了求解规模与执行稳定性。
  \item \textbf{情景与风险刻画：}对参数进行年度更新并在尾部方向加密采样，提升了 CVaR 估计的稳健性，使“期望—尾部”的权衡更加可控。
  \item \textbf{与鲁棒优化的关系：}CVaR 以概率加权方式约束尾部，鲁棒预算 \(\Gamma\) 以内集最坏情景给出确定性保证，二者在稳健性目标一致、保守度与可解释性各具优势。
  \item \textbf{数值发现：}在相近期望收益下，滚动 + CVaR 显著改善最坏 5\% 分位与 CVaR 指标，并降低对价格弹性与容量参数的尾部敏感度。
  \item \textbf{实施建议：}将 “CVaR@95\%/最坏分位” 等尾部指标纳入年度考核，同时设定地块切换频率上限，以兼顾收益与实施稳定性。
\end{itemize}

综合而言，滚动 + CVaR 框架为跨年决策与资源配置提供了可操作的风险控制工具，也为不同风险偏好的管理者提供了透明的调参把手（如 \(\alpha,\lambda,\kappa\) 与分段点密度）。

% ==================== 本章结束 ====================